% Pursuing Stacks (cotes 134-2 to 134-8): archive cote 134-4, site page 38 (PDF page 39), typescript page 126.
% typescript. Draft compared: ps3.tex:1891-1967.
% AI-transcribed by Claude Opus 5.5 (claude-opus-5-5, adaptive thinking, effort high), Message Batches API, batch ps2 (msgbatch_01NVwr1Yuo8qJf5rykeGBXPK), 2026-11-12; not yet checked by eye.
\scanpage{134-4/38}
\mspage{126}
a)\quad T\over{H}{h}e functor M\hand{$\to$}M' is bijective on objects,  surjective on arrows.
b)\quad The functor M\hand{$\to$}M' is a localization functor \hand{(}it will be so then,  in
view of a),  in the strict sense,  namely M' will be \inserted{M-}isomorphic,  no\hand{t}
only \inserted{M-}equivalent, to a localization $W^{-1}M$ ).

However,  the question arises w\hand{h}ether it is possible to define \over{s}{s}uch
a homotopy structure on M,  in terms of an arbitrary R i.e. an arbitrary
quotient category having same objects  (let Q(M)=Q the set of all such R's)\hand{,}
or in terms of an arbitrary localization of M,  or what amounts to the sam
% NOTE: the right edge of the scan cuts the ends of some lines (``sam'', ``definin'', ``relatio'', ``clear'', ``injectiv'')
me,  in terms of an arbitrary strongly saturated W\hand{$\subset$}$\mathrm{Fl}$(M)  (let's call
L(M)=L  the set of all such W's).  T\hand{h}e first thing that comes to mind
here,  is that we got two natural maps
\hand{(1)} \diagram{\begin{tikzcd} Q \arrow[r, bend left, "r"] & L \arrow[l, bend left, "s"] \end{tikzcd}}
between Q and L,  which are defined by the observation that whenever we
have a functor \inserted{i:} M\hand{$\to$}M' , \inserted[typed]{injective on objects,}\hand{,} it defines both an R\hand{$\in$}Q  (namely f\hand{$\sim$}$_{R}$g  iff
i(f)=i(g))  and a W\hand{$\in$}L (namely f\hand{$\in$}W  iff u(f) is invertible).  \over{f}{F}or definin
r(R)  resp. s(\over{\illegible{1}}{W}),  we apply this to the case when M' = M/R  resp. $W^{-1}$M.
Let's look a little at the two cases separately.

Start with R\inserted{in Q}, we get W=r(R),
\[ W = \hand{\{} f \hand{\in} \mathrm{Fl}(M) \hand{\mid} i(f) \text{ invertible}\hand{\}} \hand{,} \]
where
\[ i:\ M \hand{\to} M/R \]
is the canonical functor;  we thus get a canonical functor (compatible
with the structures ``under M'' )
\[ (\over{*}{2}) \qquad \deleted{\xout{\illegible{1}}}\hand{$\alpha=\alpha_{R}$}:\ M_W \hand{\to} M_R \hand{,} \]
where for simplicity I write
\[ M_R = M/R \ ,\ \ M_W = W^{-1}M \ . \]
We may define W as the largest element in L (for the natural order relatio
in L namely inclusion of subsets of $\mathrm{Fl}$(M))  such that a functor (*)
exists (compatible with the functors from M into both sides)  -  such a
functor of course is unique (by the preliminaries on functors M\hand{$\to$}M' made
at the beginning).  \xout{Now} In terms of (\over{\unsure{*}}{1}),  we can say that R is \inserted{actually} a homotopy
relation  (let's call $Q_o$(M)=$Q_o$ the \xout{corresponding} subset of Q\deleted{\unsure{'s}} of all such
\xout{Now start with}
relations)  iff \deleted{(*} (\hand{2}) is an isomorphism, or equivalently (as it is clear
ly bijective on \xout{\illegible{1}} objects,  sur\hand{j}ective on arrows)  if\hand{f} it is \uline{injectiv
on arrows}\hand{,} i.e. \uline{faithful}.

Conversely,  start now with W in L ,  we get R=s(W) ,
\[ R = \hand{\{}(f,g)\hand{\in} \mathrm{Fl}(M)\mathrm{x}\mathrm{Fl}(M) \hand{\mid} i(f)=i(g)\hand{\}} \ , \]
\margin{38}
